\documentclass[11pt]{article}
\usepackage{arxiv}
\usepackage[utf8]{inputenc}
\usepackage{amsmath,amssymb,amsfonts}
\usepackage{hyperref}
\usepackage{booktabs}
\usepackage{graphicx}

\title{SKI Combinator Reduction as a Cellular Automaton on Negatively Curved Graphs}
\author{Jonathan Hill\\
\texttt{qizwiz} \thanks{Independent researcher. Correspondence: qizwiz@gmail.com}}

\date{\today}

\begin{document}
\maketitle

\begin{abstract}
We introduce a cellular automaton designed to perform SKI combinator reduction on graph-structured cell spaces. Our main observation is a critical threshold for S-reduction, which fails on binary trees with a maximum degree of 5 and succeeds on ternary trees with a maximum degree of 7. This suggests that universal computation via combinatory logic necessitates negative curvature in the underlying cell space.
\end{abstract}

\section{Introduction}

Cellular automata (CA) are a well-studied model of computation in which a collection of cells, arranged on some fixed topology, update their states simultaneously according to local transition rules. Since Cook's proof that Rule 110 is Turing-complete [Cook 2004], it has been known that even one-dimensional CAs on flat lattices can perform universal computation. However, Cook's construction requires carefully engineered initial conditions that simulate cyclic tag systems through glider interactions -- the topology itself plays no computational role.

We ask a different question: can the \textit{topology} of a CA's cell space directly enable computation? Specifically, can the graph structure of a CA be chosen so that standard combinatory logic reduction -- the basis of functional programming -- emerges as the natural local dynamics?

SKI combinatory logic [Schonfinkel 1924, Curry 1930] provides a minimal universal basis for computation. Every computable function can be expressed using three combinators:
\begin{itemize}
\item \textbf{I} x = x (identity)
\item \textbf{K} x y = x (constant)
\item \textbf{S} x y z = (x z)(y z) (application with duplication)
\end{itemize}

I-reduction and K-reduction are local operations. S-reduction, however, requires \textit{duplicating} the subterm z, which requires an available neighboring cell -- a spatial constraint dependent on local connectivity.

\subsection{Contributions}
\begin{enumerate}
\item A CA implementation of SKI reduction on graph-structured cell spaces with double-buffered parallel updates (Section 3).
\item An empirical demonstration that S-reduction exhibits a sharp connectivity threshold: fails on binary-branching graphs (max degree 5), succeeds on ternary-branching graphs (max degree 7) due to local neighborhood surplus (Section 4).
\item A geometric interpretation relating this threshold to discrete curvature, suggesting universal computation via combinatory logic requires negative curvature (Section 5).
\end{enumerate}

\section{The SKI-CA System}

\subsection{Cell Structure}

Each cell maintains:
\begin{itemize}
\item State: \{EMPTY, S, K, I, VAR, APP\}
\item Auxiliary value (for VAR: variable name)
\item Left/right pointers (for APP: function/argument)
\item Used flag (whether cell holds active data)
\end{itemize}

All fields are double-buffered, atomically swapped after all cells compute.

\subsection{Graph Topology}

Cells arranged on regular trees with branching factor $k$, depth $d$, with sibling links:

\begin{center}
\begin{tabular}{ccc}
\toprule
$k$ & Cells (d=6) & Degree range \\
\midrule
2 & 127 & 2--5 \\
3 & 1,093 & 3--7 \\
4 & 5,461 & 3--8 \\
5 & 19,531 & 3--9 \\
\bottomrule
\end{tabular}
\end{center}

\subsection{S-Reduction Detail}

S x y z $\rightarrow$ (x z)(y z) requires:
\begin{enumerate}
\item Finding a free neighboring cell for the copy of z
\item Creating two new APP structures: (x z-original) and (y z-copy)
\item Rewiring the root to these new APPs
\end{enumerate}

Step 1 is the bottleneck. If no free neighbor exists, reduction cannot proceed.

\section{Experimental Results}

All tests via Common Lisp implementation (475 lines, 11/11 passing).

\subsection{Connectivity Threshold}

\begin{center}
\begin{tabular}{ccccc}
\toprule
$k$ & Degree range & I & K & S \\
\midrule
2 & 2--5 & PASS & PASS & \textbf{FAIL} \\
3 & 3--7 & PASS & PASS & PASS \\
4 & 3--8 & PASS & PASS & PASS \\
5 & 3--9 & PASS & PASS & PASS \\
\bottomrule
\end{tabular}
\end{center}

The threshold is sharp across all tested depths (3--6).

\subsection{Properties of the Threshold}

On branch=2 graphs after embedding S I I x: APP nodes near the redex have 4 neighbors but \textbf{0 free}. The search for a free cell for z-copy fails. On branch $\geq$ 3, the exponentially growing neighborhood provides sufficient surplus.

\section{Geometric Interpretation}

\subsection{Trees as Discrete Hyperbolic Space}

Regular trees of branching factor $k$ are discrete models of hyperbolic space [Gromov 1987]. Key property: exponential growth $O(k^r)$ vs $O(r^{d-1})$ in Euclidean $d$-space. This exponential growth is the hallmark of negative curvature.

\subsection{Curvature and Computational Capacity}

S-reduction requires local neighborhood surplus -- more cells available nearby than the term consumes. This surplus is precisely what negative curvature provides. Our empirical finding: the curvature threshold for SKI reduction lies between branch=2 and branch=3. Below this threshold, local geometry is too tight for subterm duplication. Above it, exponential neighborhood growth ensures sufficient slack.

\section{Limitations}

Our S-reduction uses shallow copying (root cell only). Deep copying complex terms remains future work. We reduce one redex per step; parallel reduction of independent redexes is also future work. The curvature threshold is empirical; formal proof that branch $\geq$ 3 is necessary and sufficient remains open.

\section{Conclusion}

We have shown that SKI combinator reduction can be implemented as a strictly local CA on graph-structured cell spaces, and that S-reduction imposes a sharp curvature requirement. Binary-branching trees lack sufficient neighborhood surplus; ternary and higher provide it. This directly connects substrate geometry to computational power.

\section*{Acknowledgments}

The author thanks the Lenia community and M. Buehler for discussions on cellular automata and computation. Implementation available upon request.

\begin{thebibliography}{9}

\bibitem{Cook2004}
Cook, M. (2004). Universality in elementary cellular automata. \textit{Complex Systems}, 15(1), 1--40.

\bibitem{Schonfinkel1924}
Sch\"{o}nfinkel, M. (1924). \"{U}ber die Bausteine der mathematischen Logik. \textit{Mathematische Annalen}, 92, 305--316.

\bibitem{Curry1930}
Curry, H.B. (1930). Grundlagen der kombinatorischen Logik. \textit{American Journal of Mathematics}, 52(3), 509--536.

\bibitem{Gromov1987}
Gromov, M. (1987). Hyperbolic groups. In \textit{Essays in Group Theory}, MSRI Publications, 8, 75--264.

\bibitem{Margenstern2007}
Margenstern, M. (2007). \textit{Cellular Automata in Hyperbolic Spaces}. OCP Science.

\bibitem{Mordvintsev2020}
Mordvintsev, A., et al. (2020). Growing Neural Cellular Automata. \textit{Distill}, 5(2).

\bibitem{Buehler2024}
Buehler, M. (2024). LifeGPT: Language modeling as a life simulator. arXiv:2409.12182.

\bibitem{Randazzo2024}
Randazzo, E., et al. (2024). Computational Life: self-replicators emerging on computational substrates. arXiv:2406.19108.

\bibitem{Lu2024}
Lu, C., et al. (2024). JaxLife: Evolved agents in Turing-complete environments. arXiv:2409.00853.

\end{thebibliography}

\end{document}
